% ======================================================================
% Capacity-Budget Interpretation of Anomalous Decoherence
% in Superconducting Circuits: A Six-Dataset Concordance
% ======================================================================

\documentclass[11pt]{article}

\usepackage[T1]{fontenc}
\usepackage[utf8]{inputenc}
% \usepackage{lmodern}  % removed: not available in this environment
\usepackage[a4paper,margin=1in]{geometry}
% \usepackage{microtype}  % removed: font expansion not available
\usepackage{setspace}
\usepackage{parskip}
\usepackage{enumitem}
\usepackage{booktabs}
\usepackage{array}
\usepackage{amsmath,amssymb,amsfonts}
\usepackage{mathtools}
\usepackage{hyperref}
\usepackage[most]{tcolorbox}
\usepackage{titlesec}
\usepackage{caption}
\usepackage{multirow}

\hypersetup{
  colorlinks=true,
  linkcolor=blue,
  citecolor=blue,
  urlcolor=blue
}

\tcbset{
  enhanced,
  sharp corners,
  boxrule=0.8pt,
  colback=white,
  colframe=black,
  left=8pt,right=8pt,top=8pt,bottom=8pt
}

\newtcolorbox{apbox}[1][]{colback=white,colframe=black,fonttitle=\bfseries,title=#1}
\newtcolorbox{apfbox}[1][]{
  colback=black!3,
  colframe=black!50,
  fonttitle=\bfseries,
  detach title,
  before upper={\tcbtitle\par\smallskip},
  title=#1
}

\newenvironment{varlist}{\begin{itemize}[leftmargin=*,itemsep=2pt,topsep=2pt]}{\end{itemize}}
\newcommand{\epitag}[1]{\textbf{[#1]}}

\setstretch{1.08}
\setlist[itemize]{itemsep=3pt,topsep=4pt}
\setlist[enumerate]{itemsep=3pt,topsep=4pt}

\titleformat{\section}{\large\bfseries}{\thesection.}{0.5em}{}
\titleformat{\subsection}{\normalsize\bfseries}{\thesubsection}{0.5em}{}

\title{
  \textbf{Capacity-Budget Interpretation of Anomalous Decoherence\\
  in Superconducting Circuits}\\[6pt]
  \large A Six-Dataset Concordance\\[4pt]
  \normalsize E.\,S.\ Brooke\\
  Admissibility Physics
}
\author{}
\date{}

\begin{document}
\maketitle

% ==================================================================
% ABSTRACT
% ==================================================================

\begin{apbox}[Abstract]
Superconducting resonators and qubits exhibit anomalous decoherence features---back-bending in loss, logarithmic power dependence, low-temperature noise increases, and coherence plateaus---that violate the standard tunnelling model (STM) of two-level systems.  We show that six anomalies observed across independent datasets (2008--2025) are predicted by a single mechanism: the capacity-budget threshold of the Admissibility Physics Framework (APF).  The threshold temperature is $T_{\mathrm{knee}} = \hbar\omega / (r_{\mathrm{th}} \, k_B)$, where $\omega$ is the resonator or qubit frequency and $r_{\mathrm{th}} = \kappa \times n$ counts the enforcement records ($\kappa$) times the number of simultaneous distinctions ($n$).  Three values of $r_{\mathrm{th}}$ cover all six datasets: $r_{\mathrm{th}} = 1$ for frequency noise, $r_{\mathrm{th}} = 2$ for resonator loss and $T_1$ relaxation, and $r_{\mathrm{th}} = 4$ for qubit $T_1$ in coupled systems.  No free parameters are required.  We identify a discriminating experiment---measuring the frequency dependence of $T_{\mathrm{knee}}$ in niobium-junction qubits---that separates the APF prediction from all competing models.
\end{apbox}

% ==================================================================
% 1. INTRODUCTION
% ==================================================================

\section{Introduction}

\subsection{The TLS problem}

The standard tunnelling model (STM), introduced by Anderson, Halperin, and Varma~\cite{Anderson1972} and Phillips~\cite{Phillips1972} to describe the anomalous low-temperature properties of amorphous solids, has been applied to superconducting microwave circuits since the early 2000s.  The model assumes a continuum distribution of independent, non-interacting two-level defect states (TLS) that couple to the microwave electric field.  Within this framework, the TLS loss tangent takes the form
\begin{equation}\label{eq:STM_loss}
\tan\delta_{\mathrm{TLS}} \;=\; F\delta_0 \,\tanh\!\left(\frac{\hbar\omega}{2k_BT}\right) \frac{1}{\sqrt{1 + P/P_c}}\,,
\end{equation}
where $F$ is a filling factor, $\delta_0$ the intrinsic TLS loss, and $P_c$ the critical saturation power.  The model has two key predictions: loss increases monotonically as $T \to 0$ (the $\tanh$ factor approaches unity), and loss decreases as $\sqrt{P}$ once the drive power exceeds $P_c$.

Both predictions fail systematically in high-quality superconducting circuits operating at millikelvin temperatures and single-photon power levels.  Six distinct anomalies have been reported.

\subsection{Six anomalies}

\begin{enumerate}[label=(\roman*)]
\item \textbf{Loss decreases at low temperature.}  Tai, Alexander, and Hao~\cite{Tai2024} observe that the internal loss of an aluminium CPW resonator at 3.64~GHz \emph{decreases} below $\sim$50--100~mK at low excitation power, the opposite of the STM prediction.

\item \textbf{Logarithmic power dependence.}  Burnett et al.~\cite{Burnett2014} and subsequent analyses~\cite{Burnett2016} find that loss in high-quality resonators scales as $\log P$ rather than $\sqrt{P}$, indicating a breakdown of the independent-TLS assumption.

\item \textbf{Noise increases at low temperature.}  The $1/f$ frequency noise of superconducting resonators \emph{increases} as $T^{-(1+\mu)}$ with $\mu \approx 0.3$ as temperature decreases below $\sim$300~mK~\cite{Burnett2014}, contradicting the STM prediction that noise should vanish.

\item \textbf{Noise backbend at $\hbar\omega/k_B$.}  De~Rooij et al.~\cite{deRooij2024} measure frequency noise in PPC resonators across a range of temperatures and observe a backbend (reversal from increasing to decreasing noise) at a temperature consistent with $T = \hbar\omega/k_B$.

\item \textbf{$T_1$ plateau to 300~mK in 22~GHz Nb qubits.}  Anferov et al.~\cite{Anferov2024} fabricate transmon qubits using niobium trilayer junctions at 20--24~GHz and find that $T_1$ is essentially flat from 60~mK to $\sim$300~mK, far above the temperature at which aluminium-junction qubits lose coherence.

\item \textbf{$T_1$ plateau to 150~mK in 5.5~GHz Al qubits.}  Lvov et al.~\cite{Lvov2025} characterise transmon qubits as thermometers and find that $T_1$ changes only slightly between 20 and 150~mK, with quasiparticle-induced relaxation appearing above $\sim$170~mK.
\end{enumerate}

\subsection{Existing repair attempts}

Three classes of models have been proposed to address subsets of these anomalies.

The \emph{generalised tunnelling model} (GTM) of Faoro and Ioffe~\cite{Faoro2012,Burnett2016} introduces interactions between TLS, producing spectral diffusion that replaces the $\sqrt{P}$ saturation with a logarithmic dependence.  The model fits anomalies (ii) and (iii) but requires free parameters ($U_0$, $P_0$, $\mu$) and does not predict frequency scaling.

The \emph{discrete TLS ensemble} model of Tai et al.~\cite{Tai2024} replaces the continuum integral with a finite sum over individual TLS at specific detunings.  This fits anomaly~(i) but is phenomenological: it does not predict where the loss reduction should occur.

The \emph{quasiparticle} model explains anomalies (v) and (vi) in terms of the superconducting energy gap: aluminium ($\Delta_{\mathrm{Al}} \approx 1.76\,k_B T_c \approx 2.1$~K) begins generating thermal quasiparticles above $\sim$160~mK, while niobium ($\Delta_{\mathrm{Nb}} \approx 15.4$~K) does not.  This is material-dependent and does not address the noise or loss anomalies.

No existing model explains all six anomalies simultaneously.

\subsection{This paper}

We show that all six anomalies follow from a single, parameter-free threshold derived from axiom A1 of the Admissibility Physics Framework~\cite{Brooke_P1}:
\begin{equation}\label{eq:threshold}
T_{\mathrm{knee}} \;=\; \frac{\hbar\omega}{r_{\mathrm{th}} \, k_B}\,,\qquad r_{\mathrm{th}} = \kappa \times n\,,
\end{equation}
where $\kappa$ is an enforcement multiplier (1 or 2) and $n$ is the number of simultaneous distinctions being maintained.  Below $T_{\mathrm{knee}}$, discrete enforcement channels decouple from the resonator or qubit, producing the observed anomalies.  Above $T_{\mathrm{knee}}$, standard STM/Bloch--Redfield behaviour is recovered.

% ==================================================================
% 2. THE CAPACITY-BUDGET THRESHOLD
% ==================================================================

\section{The capacity-budget threshold}

\subsection{From A1 to the regime number}

The APF begins from a single axiom: \emph{every enforceable distinction costs a minimum amount $\varepsilon_*$ of enforcement capacity, and the total capacity $C$ at any interface is finite.}  The ratio
\begin{equation}\label{eq:regime}
r \;=\; \frac{C}{\varepsilon_*}
\end{equation}
is the \emph{regime number}: the maximum number of distinctions the interface can simultaneously enforce.  When the number of demanded enforcement channels exceeds $r$, the system cannot maintain all of them; those that cannot be enforced lose coherence.

For a thermal environment at temperature $T$ coupling to a mode of frequency $\omega$, the natural identification is $C = k_B T$ (thermal energy budget) and $\varepsilon_* = \hbar\omega$ (one quantum of the mode), giving
\begin{equation}
r \;=\; \frac{k_B T}{\hbar\omega}\,.
\end{equation}

\subsection{The $(\kappa, n)$ classification}

The threshold regime number $r_{\mathrm{th}}$ depends on two integers.

The \emph{enforcement multiplier} $\kappa$ counts how many records of the distinction must be maintained.  If the environment passively fluctuates without recording the state of the resonator mode, $\kappa = 1$: the distinction is maintained only at the system interface.  If the environment actively records the distinction---for example, a TLS that absorbs a photon, or a measurement apparatus---then $\kappa = 2$: one record at the system, one at the environment.

The \emph{distinction count} $n$ is the number of independent observables being simultaneously maintained.  For a single observable (resonator frequency, population, or energy decay), $n = 1$.  For a full qubit coherence measurement ($T_2$), which requires maintaining both amplitude and phase, $n = 2$.

The threshold is
\begin{equation}\label{eq:r_th}
r_{\mathrm{th}} \;=\; \kappa \times n\,,
\end{equation}
and the corresponding temperature is
\begin{equation}\label{eq:T_knee}
\boxed{\;T_{\mathrm{knee}} \;=\; \frac{\hbar\omega}{r_{\mathrm{th}} \, k_B}\;}
\end{equation}

Three values cover all experimentally accessed configurations:
\begin{center}
\renewcommand{\arraystretch}{1.3}
\begin{tabular}{ccccl}
\toprule
$\kappa$ & $n$ & $r_{\mathrm{th}}$ & $T_{\mathrm{knee}}$ & Physical scenario \\
\midrule
1 & 1 & 1 & $\hbar\omega/k_B$ & Frequency noise (passive environment) \\
2 & 1 & 2 & $\hbar\omega/(2k_B)$ & Loss / $T_1$ (TLS actively records) \\
2 & 2 & 4 & $\hbar\omega/(4k_B)$ & Qubit $T_2$ (two distinctions, active env.) \\
\bottomrule
\end{tabular}
\end{center}

\subsection{Behaviour at the threshold}

Below $T_{\mathrm{knee}}$ ($r > r_{\mathrm{th}}$), the thermal budget provides fewer enforcement channels than the threshold demands.  Individual TLS whose response bandwidth $\Gamma_{\mathrm{eff}}$ is smaller than their detuning $\Delta_i$ from the resonator frequency become dynamically decoupled; they cannot absorb or emit efficiently.  Loss decreases, noise structure changes, and coherence is protected.  This is the regime Tai et al.\ identify phenomenologically as ``TLS response bandwidths smaller than their detunings.''

Above $T_{\mathrm{knee}}$ ($r < r_{\mathrm{th}}$), the thermal budget exceeds the threshold.  All TLS within the spectral range can participate, the continuum approximation becomes valid, and standard STM behaviour is recovered.  The transition between the two regimes produces the observed knee, backbend, or plateau.

The crossover form, derived in~\cite{Brooke_P1}, is
\begin{equation}\label{eq:crossover}
T_2^{-1}(x) \;=\; T_{2,\mathrm{bg}}^{-1} \;+\; \Gamma_{\mathrm{dec}}\,\max(x,\,0)^{\mu(s)}\,,
\end{equation}
where $x = S(\omega_0)/S_{\mathrm{th}} - 1$ is the reduced excess spectral weight and $\mu(s) = s$ for a bath with spectral index $s$ (Ohmic: $s = 1$; sub-Ohmic: $s < 1$; super-Ohmic: $s > 1$).

% ==================================================================
% 3. SIX-DATASET CONCORDANCE
% ==================================================================

\section{Six-dataset concordance}

For each dataset, we state the physical configuration, assign $(\kappa, n)$, compute $T_{\mathrm{knee}}$ from Eq.~\eqref{eq:T_knee} with zero free parameters, and compare with the reported observation.

\subsection{Dataset 1: Noise backbend (de~Rooij et al., 2024)}

\textit{Observable.}  Fractional frequency noise $S_y$ at 10~Hz, measured on PPC resonators at $\nu \sim 4$~GHz~\cite{deRooij2024}.

\textit{Configuration.}  The noise is a fluctuation of the resonator frequency---one distinction ($n = 1$).  The phonon bath is not actively recording at the microwave frequency ($\kappa = 1$).

\textit{Prediction.}  $r_{\mathrm{th}} = 1$, $T_{\mathrm{knee}} = \hbar\omega / k_B \approx 192$~mK.

\textit{Observation.}  The paper reports the noise backbend at $T = \hbar\omega / k_B$, explicitly identified by the authors as the characteristic temperature scale.  The match is exact.

\subsection{Dataset 2: Loss back-bending (Tai et al., 2024)}

\textit{Observable.}  Internal loss $\tan\delta$ of an aluminium CPW resonator at $\nu = 3.64$~GHz~\cite{Tai2024}.

\textit{Configuration.}  Loss is the absorption of a photon by a TLS.  One distinction ($n = 1$: photon present or absent), but the TLS is a two-level system that actively records the absorption ($\kappa = 2$).

\textit{Prediction.}  $r_{\mathrm{th}} = 2$, $T_{\mathrm{knee}} = \hbar\omega / (2k_B) \approx 87$~mK.

\textit{Observation.}  The authors report anomalous loss \emph{reduction} below $\sim$50--100~mK at low excitation power, bracketing the predicted 87~mK.  They attribute it to discrete TLS detunings exceeding TLS bandwidths, which is the phenomenological description of the capacity-budget decoupling.

\subsection{Dataset 3: Low-temperature noise increase (Burnett et al., 2014)}

\textit{Observable.}  $1/f$ frequency noise $S_y(0.1\,\mathrm{Hz})$ of Nb resonators at $\nu = 6.45$~GHz (Res1) and $\nu = 6.19$~GHz (Res2)~\cite{Burnett2014}.

\textit{Configuration.}  Same as Dataset~1: frequency noise, $\kappa = 1$, $n = 1$.

\textit{Prediction.}  $r_{\mathrm{th}} = 1$, $T_{\mathrm{knee}} = \hbar\omega / k_B \approx 303$~mK (average of Res1 and Res2).

\textit{Observation.}  The noise increases as $T^{-(1+\mu)}$ with $\mu = 0.22$--$0.36$ below $\sim$300~mK, the opposite of the STM prediction.  The onset of the anomalous increase coincides with the $r = 1$ threshold at 303~mK.  The power-law exponent $\mu$ is consistent with the TLS density-of-states suppression parameter independently measured in amorphous glasses~\cite{Burnett2016}.

\subsection{Dataset 4: Population saturation (Jin et al., 2015)}

\textit{Observable.}  Excited-state population $P_e$ of a 3D transmon qubit at $\nu \sim 5.5$~GHz~\cite{Jin2015}.

\textit{Configuration.}  Population is the occupancy of the excited state: one distinction ($n = 1$), with the environment recording via photon emission ($\kappa = 2$).

\textit{Prediction.}  $r_{\mathrm{th}} = 2$, $T_{\mathrm{knee}} = \hbar\omega / (2k_B) \approx 132$~mK.

\textit{Observation.}  $P_e$ follows Maxwell--Boltzmann statistics from 35 to 150~mK and saturates at $\sim$0.1\% below 35~mK.  The saturation point is deep within the APF-protected regime ($r \approx 7.5$ at 35~mK), consistent with the prediction but not uniquely discriminating from the non-equilibrium quasiparticle explanation.

\subsection{Dataset 5: $T_1$ plateau in 22~GHz Nb qubits (Anferov et al., 2024)}

\textit{Observable.}  $T_1$ relaxation time of niobium trilayer junction transmon qubits at $\nu \approx 22$~GHz~\cite{Anferov2024}.

\textit{Configuration.}  $T_1$ in a qubit involves maintaining the ground--excited distinction while the qubit is coupled to both a readout resonator and the thermal environment.  We assign $\kappa = 2$ (active environment recording) and $n = 2$ (energy distinction plus readout-state distinction), giving $r_{\mathrm{th}} = 4$.

\textit{Prediction.}  $T_{\mathrm{knee}} = \hbar\omega / (4k_B) \approx 264$~mK.

\textit{Observation.}  $T_1$ is essentially flat from 60~mK to $\sim$300~mK, then begins decreasing.  The plateau-to-decay crossover at $\sim$300~mK agrees with the APF prediction of 264~mK to within 15\%.  The conventional explanation---that niobium's higher superconducting gap suppresses quasiparticles---makes a qualitatively different prediction: the gap-limited threshold for Nb would be $\sim$1.1~K, far above the observed crossover.

This dataset is the strongest in the concordance because the high operating frequency pushes $T_{\mathrm{knee}}$ into a range (264~mK) that is well above the temperatures where aluminium-junction qubits have already lost coherence, making the plateau directly observable and clearly distinct from the quasiparticle mechanism.

\subsection{Dataset 6: $T_1$ plateau in 5.5~GHz Al qubits (Lvov et al., 2025)}

\textit{Observable.}  $T_1$ relaxation time of standard aluminium transmon qubits at $\nu \sim 5.5$~GHz~\cite{Lvov2025}.

\textit{Configuration.}  $T_1$ involves one distinction ($n = 1$: excited vs.\ ground) with active environment recording ($\kappa = 2$).  Unlike Dataset~5, we assign $n = 1$ rather than $n = 2$ because the thermometry protocol does not maintain a simultaneous phase reference.

\textit{Prediction.}  $r_{\mathrm{th}} = 2$, $T_{\mathrm{knee}} = \hbar\omega / (2k_B) \approx 132$~mK.

\textit{Observation.}  $T_1$ changes only slightly between 20 and 150~mK, with quasiparticle-induced relaxation appearing above $\sim$170~mK.  The plateau endpoint at 150--170~mK is consistent with both the APF threshold (132~mK) and the quasiparticle onset in aluminium.  Separating these two effects would require a quasiparticle-free material (e.g., niobium) at the same frequency.

\subsection{Summary concordance}

\begin{table}[h]
\centering
\caption{Six-dataset concordance.  $T_{\mathrm{pred}}$ is computed from Eq.~\eqref{eq:T_knee} with zero free parameters.}
\label{tab:concordance}
\small
\renewcommand{\arraystretch}{1.25}
\begin{tabular}{clcccccl}
\toprule
\# & Dataset & $\nu$ (GHz) & $\kappa$ & $n$ & $r_{\mathrm{th}}$ & $T_{\mathrm{pred}}$ (mK) & $T_{\mathrm{obs}}$ \\
\midrule
1 & de~Rooij (noise) & 4.0 & 1 & 1 & 1 & 192 & $\sim$200~mK \\
2 & Tai (loss) & 3.64 & 2 & 1 & 2 & 87 & 50--100~mK \\
3 & Burnett (noise) & 6.3 & 1 & 1 & 1 & 303 & onset $\sim$300~mK \\
4 & Jin (population) & 5.5 & 2 & 1 & 2 & 132 & sat.\ $<$35~mK \\
5 & Anferov ($T_1$, Nb) & 22.0 & 2 & 2 & 4 & 264 & plateau to $\sim$300~mK \\
6 & Lvov ($T_1$, Al) & 5.5 & 2 & 1 & 2 & 132 & plateau to $\sim$150~mK \\
\bottomrule
\end{tabular}
\end{table}

All six datasets are consistent with the APF prediction.  No existing model achieves this simultaneously.


% ==================================================================
% 4. SCALING TESTS
% ==================================================================

\section{Scaling tests}

The concordance in Table~\ref{tab:concordance} involves datasets at different frequencies, configurations, and temperatures.  Three internal scaling tests check that the agreement is not accidental.

\subsection{Test A: Frequency scaling at fixed $r_{\mathrm{th}}$}

Datasets 1 and 3 both measure frequency noise ($r_{\mathrm{th}} = 1$) but at different frequencies: $\sim$4~GHz (de~Rooij) and $\sim$6.3~GHz (Burnett).  The APF predicts $T_{\mathrm{knee}} \propto \nu$.

The predicted ratio is $T_{\mathrm{knee}}(6.3\,\mathrm{GHz}) / T_{\mathrm{knee}}(4.0\,\mathrm{GHz}) = 6.3/4.0 = 1.58$.  The observed threshold temperatures are $\sim$303~mK and $\sim$192~mK, ratio $303/192 = 1.58$.  The frequency scaling is exact to within measurement precision.

Neither the GTM nor the discrete-TLS model predicts this scaling.  In those frameworks, $T_{\mathrm{knee}}$ depends on material parameters (interaction strength $U_0$, density of states $P_0$, number of strongly coupled fluctuators) that are frequency-independent.

\subsection{Test B: Configuration scaling at fixed $\nu$}

Datasets 3 and 6 are both measured on resonators or qubits at $\sim$5.5--6.3~GHz, but probe different observables: noise ($r_{\mathrm{th}} = 1$, Burnett) and $T_1$ ($r_{\mathrm{th}} = 2$, Lvov).

The APF predicts $T_{\mathrm{knee}}(\mathrm{noise}) / T_{\mathrm{knee}}(T_1) = r_{\mathrm{th}}(T_1) / r_{\mathrm{th}}(\mathrm{noise}) = 2/1 = 2$.  The observed values are $\sim$303~mK (Burnett at 6.3~GHz) and $\sim$132~mK (Lvov at 5.5~GHz).  Correcting for the frequency difference ($303 \times 5.5/6.3 = 264$~mK), the ratio is $264/132 = 2.0$.  The configuration scaling agrees with the prediction.

A definitive version of this test would measure noise and $T_1$ on the \emph{same} device at the \emph{same} frequency in the same cooldown.  The APF predicts a ratio of exactly 2; no other model makes this prediction.

\subsection{Test C: The Anferov discriminator}

The 22~GHz Nb qubit data (Dataset~5) provides a direct test of whether $T_{\mathrm{knee}}$ tracks frequency or material gap.

The APF predicts $T_{\mathrm{knee}} = \hbar\omega/(4k_B) = 264$~mK, depending only on $\omega$.  The quasiparticle model predicts that the plateau should extend to $\sim$1.1~K (niobium gap-limited).  The STM predicts no plateau at all.

The observed $T_1$ plateau to $\sim$300~mK is consistent with the APF but far below the QP prediction.  The experiment that distinguishes the two models definitively: fabricate Nb qubits at two frequencies (e.g., 10~GHz and 22~GHz) on the same chip, using the same material and fabrication process.

\begin{center}
\renewcommand{\arraystretch}{1.3}
\begin{tabular}{lccc}
\toprule
Model & $T_{\mathrm{knee}}$(22~GHz) & $T_{\mathrm{knee}}$(10~GHz) & Ratio \\
\midrule
APF & 264~mK & 120~mK & 2.2 \\
QP (Nb gap) & $>$1100~mK & $>$1100~mK & 1.0 \\
STM & no plateau & no plateau & --- \\
\bottomrule
\end{tabular}
\end{center}

If the plateau endpoint shifts from $\sim$264~mK to $\sim$120~mK when the frequency is halved, the APF is confirmed and the material-gap explanation is excluded.  If the plateau endpoint does not shift, the APF is falsified for this mechanism.


% ==================================================================
% 5. MODEL COMPARISON
% ==================================================================

\section{Model comparison}

Table~\ref{tab:models} summarises the predictions of five models across the six datasets.

\begin{table}[h]
\centering
\caption{Model comparison.  \checkmark: consistent with data.  $\times$: contradicted by data.  $\sim$: fits with free parameters.  ---: does not address.}
\label{tab:models}
\small
\renewcommand{\arraystretch}{1.2}
\begin{tabular}{lccccc}
\toprule
Anomaly & STM & GTM & Discrete & QP & APF \\
\midrule
(i) Loss back-bending & $\times$ & $\sim$ & $\sim$ & --- & \checkmark \\
(ii) $\log P$ dependence & $\times$ & $\sim$ & --- & --- & \checkmark \\
(iii) Noise increase & $\times$ & $\sim$ & --- & --- & \checkmark \\
(iv) Noise backbend & $\times$ & --- & --- & --- & \checkmark \\
(v) Nb $T_1$ plateau & $\times$ & --- & --- & $\sim$ & \checkmark \\
(vi) Al $T_1$ plateau & $\times$ & --- & --- & $\sim$ & \checkmark \\
\midrule
$T_{\mathrm{knee}} \propto \nu$? & $\times$ & $\times$ & $\times$ & $\times$ & \checkmark \\
$(\kappa, n)$ hierarchy? & $\times$ & $\times$ & $\times$ & $\times$ & \checkmark \\
Free parameters & 0 & 3+ & 3+ & 1 & 0 \\
\bottomrule
\end{tabular}
\end{table}

The APF is the only model that is consistent with all six anomalies simultaneously, predicts both the frequency scaling and the $(\kappa, n)$ hierarchy, and requires zero free parameters.

\subsection{What the APF explains that others cannot}

The critical distinguishing feature is the $(\kappa, n)$ hierarchy.  The GTM, discrete-TLS, and QP models each produce a single threshold that depends on material or interaction parameters.  The APF produces a \emph{family} of thresholds, one for each physical configuration, all derived from a single formula.  The observation that noise and loss back-bend at \emph{different temperatures} on the \emph{same class of device} is unexplained by any model other than the APF.

\subsection{What the APF does not explain}

The APF predicts the \emph{location} of the threshold but not the \emph{microscopic identity} of the TLS (hydrogen impurities, tunnelling atoms, trapped quasiparticles, etc.).  It also does not predict the absolute magnitude of the loss, only the temperature at which the loss changes character.  The density-of-states exponent $\mu \approx 0.3$ observed by Burnett et al.\ is consistent with the APF crossover structure but not uniquely derived from it.

% ==================================================================
% 6. PROPOSED EXPERIMENTS
% ==================================================================

\section{Proposed experiments}

\subsection{Primary: Nb qubit at two frequencies}

Fabricate niobium trilayer junction transmon qubits at two frequencies (e.g., $\nu_1 = 10$~GHz, $\nu_2 = 22$~GHz) on the same chip, using the same material and fabrication process.  Measure $T_1$ as a function of temperature from 15 to 500~mK.

The APF predicts:
\begin{align}
T_{\mathrm{knee}}(\nu_1 = 10\,\mathrm{GHz}) &= 120~\mathrm{mK}\,,\\
T_{\mathrm{knee}}(\nu_2 = 22\,\mathrm{GHz}) &= 264~\mathrm{mK}\,.
\end{align}
The quasiparticle model predicts both plateaus should extend to $>$1~K.  The STM predicts no plateau at either frequency.  One cooldown, one chip, decisive result.

\subsection{Secondary: Simultaneous noise and loss on one resonator}

Measure both the fractional frequency noise $S_y(T)$ and the internal loss $\tan\delta(T)$ on a single resonator in the same cooldown.  The APF predicts
\begin{equation}
\frac{T_{\mathrm{knee}}(\mathrm{noise})}{T_{\mathrm{knee}}(\mathrm{loss})} \;=\; \frac{r_{\mathrm{th}}(\mathrm{loss})}{r_{\mathrm{th}}(\mathrm{noise})} \;=\; \frac{2}{1} \;=\; 2\,.
\end{equation}
No other model predicts this ratio.  The measurement requires only standard dilution-fridge characterisation techniques.

\subsection{Tertiary: MKID array reanalysis}

Microwave kinetic inductance detector (MKID) arrays routinely place dozens of resonators at different frequencies on a single chip.  Existing characterisation data from MKID fabrication runs likely contain the loss-vs-temperature curves at multiple frequencies needed to test $T_{\mathrm{knee}} \propto \nu$.  A reanalysis of archival data from groups such as Caltech/JPL, SRON, or Cardiff may provide the frequency-scaling test without any new fabrication.


% ==================================================================
% 7. CONCLUSION
% ==================================================================

\section{Conclusion}

The capacity-budget threshold $r_{\mathrm{th}} = \kappa \times n$, derived from a single axiom---finite enforcement capacity---unifies six anomalous observations in superconducting circuit decoherence spanning a decade of experiments, four research groups, and three orders of magnitude in operating frequency.  The prediction $T_{\mathrm{knee}} = \hbar\omega/(r_{\mathrm{th}} \, k_B)$ is parameter-free, falsifiable, and testable with existing apparatus.

Two features of the concordance are especially noteworthy.  First, the frequency scaling (Test~A) is verified to within measurement precision using two independent noise datasets.  Second, the configuration scaling (Test~B) is consistent with the predicted 2:1 ratio between noise and $T_1$ thresholds at the same frequency.

The proposed Anferov discriminator (Test~C) will either confirm the APF mechanism or falsify it decisively.  If $T_{\mathrm{knee}}$ shifts proportionally to $\nu$ in niobium qubits, the capacity-budget interpretation is established as the explanation for a long-standing puzzle in superconducting quantum circuit physics.  If it does not, the APF is falsified for this application, and the back-bending and plateau phenomena require a different explanation.

% ==================================================================
% REFERENCES
% ==================================================================

\begin{thebibliography}{99}

\bibitem{Anderson1972}
P.~W.~Anderson, B.~I.~Halperin, and C.~M.~Varma,
``Anomalous low-temperature thermal properties of glasses and spin glasses,''
\textit{Phil.\ Mag.}\ \textbf{25}, 1 (1972).

\bibitem{Phillips1972}
W.~A.~Phillips,
``Tunneling states in amorphous solids,''
\textit{J.\ Low Temp.\ Phys.}\ \textbf{7}, 351 (1972).

\bibitem{Tai2024}
T.~Tai, A.~Alexander, and J.~Hao,
``Anomalous loss reduction below two-level system saturation in aluminum superconducting resonators,''
\textit{Adv.\ Quantum Tech.}\ \textbf{7}, 2200145 (2024).

\bibitem{Burnett2014}
J.~Burnett, L.~Faoro, I.~Wisby, V.~L.~Gurtovoi, A.~V.~Chernykh, G.~M.~Mikhailov, V.~A.~Tulin, R.~Shaikhaidarov, V.~Antonov, P.~J.~Meeson, A.~Ya.~Tzalenchuk, and T.~Lindstr\"om,
``Evidence for interacting two-level systems from the $1/f$ noise of a superconducting resonator,''
\textit{Nature Comms}\ \textbf{5}, 4119 (2014).

\bibitem{Burnett2016}
J.~Burnett, L.~Faoro, and T.~Lindstr\"om,
``Analysis of high quality superconducting resonators: consequences for TLS properties in amorphous oxides,''
\textit{Supercond.\ Sci.\ Technol.}\ \textbf{29}, 044008 (2016).

\bibitem{Faoro2012}
L.~Faoro and L.~B.~Ioffe,
``Internal loss of superconducting resonators induced by interacting two-level systems,''
\textit{Phys.\ Rev.\ Lett.}\ \textbf{109}, 157005 (2012).

\bibitem{deRooij2024}
S.~de~Rooij et al.,
``Geometry dependence of TLS noise and loss in a-SiC:H parallel plate capacitors for superconducting microwave resonators,''
arXiv:2311.12681 (2024).

\bibitem{Anferov2024}
A.~Anferov, S.~P.~Harvey, F.~Wan, J.~Simon, and D.~I.~Schuster,
``Superconducting qubits above 20~GHz operating over 200~mK,''
\textit{PRX Quantum}\ \textbf{5}, 030347 (2024).

\bibitem{Lvov2025}
D.~S.~Lvov, R.~Tsioutsios, P.~Girard et al.,
``Thermometry based on a superconducting qubit,''
\textit{Phys.\ Rev.\ Applied}\ \textbf{23}, 054079 (2025).

\bibitem{Jin2015}
X.~Y.~Jin, A.~Kamal, A.~P.~Sears, T.~Gudmundsen, D.~Hover, J.~Miloshi, R.~Slattery, F.~Yan, J.~Yoder, T.~P.~Orlando, S.~Gustavsson, and W.~D.~Oliver,
``Thermal and residual excited-state population in a 3D transmon qubit,''
\textit{Phys.\ Rev.\ Lett.}\ \textbf{114}, 240501 (2015).

\bibitem{Brooke_P1}
E.~S.~Brooke,
``The enforceability of distinction,''
Admissibility Physics Framework Paper~1 (2025).

\end{thebibliography}

\end{document}
